\documentclass{situationist-preprint}

\title{The Scope of Contingency}
\msid{TSM-11-a}
\msfamily{TSM-11}
\msdiscipline{Time, Space and Modality}
\mstype{Principal manuscript}
\mskeywords{contingency; necessity; essence; necessitarianism; laws of nature; recombination}

\begin{document}
\maketitle

\begin{abstract}
Could things have been otherwise? Necessitarians, from Spinoza to recent defenders of modal collapse, answer that they could not. We argue that most things could have been otherwise, and that the necessary truths are those of logic, of mathematics and of essence. The argument combines the account of necessity of TSM-10, on which a proposition is necessary exactly when it follows logically from the essential truths, with the sparseness of essence: the real definition of a thing states what it is, and what a thing is settles neither its location, its accidental qualities, its relations to most other things, nor the laws it obeys. A principle of independence for what essences leave open then yields that every truth not derivable from the essential truths has a possible negation. We answer the objections from the principle of sufficient reason and from the essence of the cosmos.
\end{abstract}

\section{The question}

Spinoza held that things could not have been produced by God in any other way or order than that in which they were produced \citep[I, prop.~33]{Spinoza1677}. Against this stands the conviction that it might have rained yesterday and that the constants of nature might have had other values. We vindicate the conviction and say where contingency ends.

\section{The argument}

We assume the account of TSM-10: a proposition is metaphysically necessary if and only if it follows logically from the essential truths $E$, the propositions true in virtue of the natures of some objects \citep{Fine1994}.

\begin{assumption}[Sparseness]\label{a:sparse}
The real definition of a concrete thing states what it is: its kind, its constitutive origin and the like. It does not include its location, its accidental qualities, its relations to things it does not involve in its definition, or the laws of nature.
\end{assumption}

\begin{assumption}[Independence]\label{a:indep}
If a proposition about the location, accidental qualities or contingent relations of concrete things does not follow from $E$, its negation is logically consistent with $E$.
\end{assumption}

Assumption~\ref{a:sparse} is the view of MET-13 that real definitions state what things are and do not individuate them.

\begin{theorem}\label{t:main}
\textup{(i)} The necessary truths are exactly the truths of logic, the truths of mathematics, the essential truths and their logical consequences. \textup{(ii)} Every truth about the location, accidental qualities or contingent relations of concrete things that does not follow from $E$ is contingent.
\end{theorem}

\begin{proof}
(i) By TSM-10, the necessary truths are the logical consequences of $E$. Logical truths follow from any premises; mathematical truths follow from the real definitions of mathematical objects, which are in $E$, or from logic alone; essential truths are in $E$. Conversely, a consequence of $E$ is a consequence of logic, mathematics and essence. (ii) Let $p$ be such a truth. By Assumption~\ref{a:indep}, $\neg p$ is consistent with $E$, so by TSM-10 $\neg p$ is possible and $p$ is contingent.
\end{proof}

\begin{corollary}\label{c:most}
Most truths are contingent.
\end{corollary}

\begin{proof}
A concrete thing's real definition fixes a few truths, its kind, its origin and what follows; by Assumption~\ref{a:sparse} it leaves open its location, its accidental qualities and its relations to the many things it does not involve. Each of these truths is contingent by Theorem~\ref{t:main}(ii), and they outnumber the necessary ones.
\end{proof}

The laws of nature are contingent: MET-12 takes them to be the axioms of the best system summarizing the mosaic of events, which no real definition fixes. The result is neutral between contingentism and necessitism, as is MET-01.

\section{Objections}

\begin{objection}[Sufficient reason]
If every truth has a sufficient reason, and the reason for every contingent truth is ultimately a necessary truth, then every truth follows from necessary truths and is necessary.
\end{objection}

\begin{reply}
We reject the first premise in the strong form required: a contingent truth may be explained by contingent truths and laws, without a necessary truth that entails it. MET-01 gives an explanation of why there is something that is not of the kind the principle demands.
\end{reply}

\begin{objection}[The essence of the cosmos]
The cosmos is an object and has a real definition. If its definition is \enquote{the totality of what there is, arranged as it is}, everything true of it is essential, and the collapse returns.
\end{objection}

\begin{reply}
By Assumption~\ref{a:sparse}, the real definition of the cosmos states what it is, a totality of concrete things, not how they are arranged; to include the arrangement is to assume the conclusion at issue. MET-13 answers the objection from modal collapse in the same way.
\end{reply}

\section{Conclusion}

Things could have been otherwise in most respects: in where they are, how they are qualified, how they are related and which laws they obey. What could not have been otherwise is what follows from logic, from mathematics and from what things are.

\begin{thebibliography}{9}\small
\bibitem[Fine(1994)]{Fine1994} Fine, K. (1994). Essence and modality. \emph{Philosophical Perspectives}, 8, 1--16.
\bibitem[Spinoza(1677)]{Spinoza1677} Spinoza, B. (1677). \emph{Ethica Ordine Geometrico Demonstrata}. In \emph{Opera Posthuma}.
\end{thebibliography}

\end{document}
